\pagestyle{empty}
\tight
\begin{tblr}{width=\textwidth,colspec={|Q[c,m,1.9cm]|X[l,m]|},hlines,vlines,hspan=minimal,row{1}={bg=headblue},rowsep=1pt,colsep=3pt}
\SetCell{c,m}\hfil\logo\hfil & \textbf{POLITEKNIK NEGERI MALANG}\par\textbf{JURUSAN TEKNOLOGI INFORMASI}\par\textbf{PROGRAM STUDI: D4 TEKNIK INFORMATIKA} \\
\SetCell[c=2]{c} \textbf{RENCANA PEMBELAJARAN SEMESTER (RPS)} & \\
\end{tblr}
\begin{longtblr}{width=\textwidth,colspec={|Q[l,m,1.9cm]|X[0.85,l,m]|X[1.45,l,m]|X[0.75,l,m]|X[1.15,l,m]|},hlines,vlines,hspan=minimal,rowsep=1pt,colsep=2pt,row{1,3}={bg=headgray,font=\bfseries}}
MATA KULIAH & KODE & BOBOT (SKS)/jam & SEMESTER & TGL. PENYUSUNAN \\
Rekayasa Perangkat Lunak & RTI252005 & 2 SKS/4 jam & 2 & 2 Februari 2026 \\
OTORISASI & Dosen Pengembang RPS & Koordinator RMK & \SetCell[c=2]{l} Ka PRODI &  \\
 & Eka Larasati Amalia, S.ST., M.T.\par Rokhimatul Wakhidah, S.Pd., M.T.\par Vivi Nur Wijayaningrum, S.Kom., M.Kom.\par Candra Bella Vista, S.Kom., M.T.\par Dr. Eng. Banni Satria Andoko, S.Kom., M.MSI. &  & \SetCell[c=2]{l}{\vspace{8mm}\par Dr. Ely Setyo Astuti, S.T., M.T.} &  \\
\SetCell[r=16]{l} Capaian Pembelajaran (CP) & \SetCell[c=4]{l,bg=headgray,font=\bfseries} Capaian Pembelajaran Lulusan Program Studi (CPL-Prodi) &  &  &  \\
 & CPL02 & \SetCell[c=3]{l} Mampu menerapkan prinsip rekayasa sistem dan perangkat lunak untuk menghasilkan solusi aplikasi multi-platform sesuai kebutuhan. &  &  \\
 & CPL06 & \SetCell[c=3]{l} Mampu merancang, mengimplementasikan, menguji, melakukan deployment, memelihara, serta menjamin mutu perangkat lunak yang menjawab permasalahan dan memenuhi kebutuhan pengguna. &  &  \\
 & CPL08 & \SetCell[c=3]{l} Mampu mengelola proyek teknologi informasi secara profesional melalui perencanaan, koordinasi, komunikasi, dan optimalisasi sumber daya. &  &  \\
 & \SetCell[c=4]{l,bg=headgray,font=\bfseries} Capaian Pembelajaran mata kuliah (CPL-MK) &  &  &  \\
 & CPMK0205 & \SetCell[c=3]{l} Mampu menerapkan prinsip dan pola desain perangkat lunak untuk membuat blueprint arsitektural dan komponen solusi aplikasi multi-platform yang modular, dapat dipelihara (maintainable), dan dapat diperluas (scalable). &  &  \\
 & CPMK0211 & \SetCell[c=3]{l} Mampu menerapkan metodologi analisis dan desain sistem secara terstruktur untuk mentransformasi kebutuhan stakeholder menjadi spesifikasi teknis yang komprehensif bagi pengembangan aplikasi multi-platform, dengan menghasilkan artefak desain yang mencakup aspek fungsional, non-fungsional, dan arsitektural yang dapat diimplementasikan. &  &  \\
 & CPMK0602 & \SetCell[c=3]{l} Mampu merancang arsitektur dan komponen perangkat lunak secara sistematis menggunakan pola desain yang sesuai dengan kebutuhan stakeholder. &  &  \\
 & CPMK0607 & \SetCell[c=3]{l} Mampu menganalisis kebutuhan organisasi untuk merancang serta mengimplementasikan solusi sistem informasi terintegrasi yang kompleks. &  &  \\
 & CPMK0802 & \SetCell[c=3]{l} Mampu mengelola proses rekayasa perangkat lunak dan desain antarmuka (UI/UX) dalam siklus hidup proyek, termasuk penerapan metodologi pengembangan dan prinsip penjaminan mutu (SQA) untuk menghasilkan produk yang berkualitas dan berpusat pada pengguna. &  &  \\
 & \SetCell[c=4]{l,bg=headgray,font=\bfseries} Kemampuan akhir tiap tahapan belajar (Sub-CPMK) &  &  &  \\
 & SCPMK0802-01401 & \SetCell[c=3]{l} Mahasiswa mampu menjelaskan prinsip dasar manajemen proses RPL dan tahapan SDLC serta mengelola proses pengembangan melalui metodologi preskriptif (Prototyping/Spiral/RUP), metodologi Agile (Scrum/Kanban) dengan manajemen alur kerja visual, penjaminan mutu (SQA) melalui pengujian, dan pemeliharaan/evolusi sistem pasca-rilis. &  &  \\
 & SCPMK0607-01402 & \SetCell[c=3]{l} Mahasiswa mampu menganalisis kebutuhan organisasi dan melakukan elisitasi untuk menghasilkan spesifikasi kebutuhan perangkat lunak (SRS). &  &  \\
 & SCPMK0211-01403 & \SetCell[c=3]{l} Mahasiswa mampu mentransformasi kebutuhan menjadi desain teknis menggunakan pemodelan sistem (UML/Diagram). &  &  \\
 & SCPMK0602-01404 & \SetCell[c=3]{l} Mahasiswa mampu merancang arsitektur sistem dan memilih pola desain (Design Patterns) yang sistematis. &  &  \\
 & SCPMK0205-01405 & \SetCell[c=3]{l} Mahasiswa mampu merealisasikan blueprint menjadi komponen aplikasi yang modular, maintainable, dan sesuai standar kode. &  &  \\
\SetCell[r=2]{l} Penilaian & \SetCell[c=4]{l} *Nilai CPMK didapat dari agregasi nilai SUB-CPMK yang dibebankan &  &  &  \\
 & \SetCell[c=4]{l}{\tiny\begin{tblr}{width=\linewidth,colspec={|Q[c,m,0.1\linewidth]|Q[c,m,0.16\linewidth]|X[l,m]|X[l,m]|X[l,m]|X[l,m]|X[l,m]|X[l,m]|Q[c,m,0.08\linewidth]|},hlines,vlines,hspan=minimal,rowsep=1pt,colsep=0.7pt,row{1,2}={bg=headgray,font=\bfseries}}
Kode CPMK & Kode SCPMK & \SetCell[c=6]{c} Bobot per Bentuk Penilaian & & & & & & Total Bobot \\
 & & Tugas & Kuis & UTS & Case Method & PBL & UAS & \\
CPMK0205 & SCPMK0205-01405 & & & & & 20\%: implementasi kode modular dan clean code & & 20\% \\
CPMK0211 & SCPMK0211-01403 & 5\%: dokumen desain sistem (SDD/UML) & & 5\%: transformasi kebutuhan menjadi diagram UML & & & & 10\% \\
CPMK0602 & SCPMK0602-01404 & & & & & 10\%: arsitektur dan design pattern & 10\%: penerapan arsitektur pada produk akhir & 20\% \\
CPMK0607 & SCPMK0607-01402 & 10\%: dokumen SRS & & & & & & 10\% \\
CPMK0802 & SCPMK0802-01401 & 5\%: laporan pengujian & 2.5\%: konsep RPL, SDLC, metodologi & 2.5\%: pemilihan metodologi & 5\%: analisis model pengembangan & 15\%: Scrum dan Kanban & 10\%: kualitas proses dan dokumentasi & 40\% \\
\SetCell[c=8]{r}\textbf{Total Per Penilaian} & & & & & & & & \textbf{100\%} \\
\end{tblr}} &  &  &  \\
Diskripsi Singkat Mata Kuliah & \SetCell[c=4]{l} Mata kuliah ini membahas proses pengembangan perangkat lunak secara sistematis, mulai dari analisis kebutuhan, perancangan, implementasi, pengujian, hingga pemeliharaan. Materi mencakup model pengembangan seperti Waterfall, Prototype, dan Agile, serta pemodelan dengan DFD, UML, dan flowchart. &  &  &  \\
Bahan Kajian / Materi Pembelajaran & \SetCell[c=4]{l}\begin{enumerate}\item Pengantar RPL dan Software Development Life Cycle (SDLC)\item Waterfall, Prototype, Spiral, dan RUP\item Agile, Scrum, dan Kanban\item Pseudocode dan Flowchart\item Data Flow Diagram dan Context Diagram\item UML: Use Case, Use Case Specification, Activity Diagram, dan Sequence Diagram\item Kebutuhan Perangkat Lunak\item Desain Perangkat Lunak\item Implementasi\item Testing\item Evolusi Perangkat Lunak\end{enumerate} &  &  &  \\
\SetCell[r=2]{l} Pustaka & \SetCell[c=4]{l} \textbf{Utama:}\begin{enumerate}\item Ian Sommerville. 2015. Software Engineering, 10th Edition. Pearson.\item William R. King. 2015. Planning for Information Systems. Routledge.\end{enumerate} &  &  &  \\
 & \SetCell[c=4]{l} \textbf{Pendukung:}\begin{enumerate}\item Sprague, R.H. and McNurlin, B.C. 2002. Information Systems Management in Practice, 5th edition. Prentice-Hall.\item Ward, J. et al. 1996. Strategic Planning for Information Systems Practice, 2nd edition. Wiley.\end{enumerate} &  &  &  \\
Media Pembelajaran & \SetCell[c=2]{l} \textbf{Software:} StarUML/Draw.io, Trello/Jira/GitHub Projects, IDE, Git, dan LMS. &  & \SetCell[c=2]{l} \textbf{Hardware:} Komputer, laptop, dan proyektor. &  \\
Nama Dosen Pengampu & \SetCell[c=4]{l}\begin{enumerate}\item Eka Larasati Amalia, S.ST., M.T.\item Rokhimatul Wakhidah, S.Pd., M.T.\item Vivi Nur Wijayaningrum, S.Kom., M.Kom.\item Candra Bella Vista, S.Kom., M.T.\item Dr. Eng. Banni Satria Andoko, S.Kom., M.MSI.\end{enumerate} &  &  &  \\
Matakuliah Syarat & \SetCell[c=4]{l} - &  &  &  \\
\end{longtblr}
\begin{longtblr}{width=\textwidth,colspec={|Q[c,m,0.06\textwidth]|X[1.35,l,m]|X[1.08,l,m]|X[1.16,l,m]|X[0.7,l,m]|X[1.24,l,m]|X[1.06,l,m]|X[0.98,l,m]|Q[c,m,0.05\textwidth]|},hlines,vlines,hspan=minimal,rowhead=2,row{1,2}={bg=headgreen,font=\bfseries},rowsep=1pt,colsep=1.5pt}
Mgg.\par Ke & Kemampuan Akhir Yang Direncanakan (Sub-CP-MK) & Bahan kajian (Materi Pembelajaran) & Bentuk dan Metode Pembelajaran & Estimasi Waktu & Pengalaman Belajar Mahasiswa & Kriteria \& Bentuk Penilaian & Indikator Penilaian & Bobot Penilaian (\%) \\
(1) & (2) & (3) & (4) & (5) & (6) & (7) & (8) & (9) \\
1 & SCPMK0802-01401 & Pengantar RPL dan konsep dasar Software Development Life Cycle (SDLC). & Bentuk: Ceramah, diskusi, studi kasus, dan latihan analisis. Metode: Contextual Instruction (CI). & 1 x (4 x 50') tatap muka; tugas terstruktur dan mandiri. & Mahasiswa mengikuti ceramah dan diskusi tentang konsep RPL, aktif bertanya, dan mencatat poin penting. & Kriteria: ketepatan dan penguasaan. Bentuk: keaktifan diskusi kelompok dan ketepatan jawaban tugas. & Keaktifan bertanya; kelengkapan catatan. & 0 \\
2 & SCPMK0802-01401 & Tahapan SDLC dan analisis studi kasus proyek perangkat lunak nyata (sukses/gagal). & Bentuk: Ceramah, diskusi, studi kasus, dan latihan analisis. Metode: Contextual Instruction (CI). Penugasan: analisis model SDLC pada studi kasus nyata dan faktor keberhasilan/kegagalannya. & 1 x (4 x 50') tatap muka; tugas terstruktur dan mandiri. & Mahasiswa mengkaji SDLC, mengaitkan model ke kasus nyata, dan berdiskusi kelompok. & Kriteria: ketepatan dan penguasaan. Bentuk: keaktifan diskusi kelompok dan ketepatan jawaban tugas. & Ketepatan identifikasi model SDLC dan faktor keberhasilan/kegagalan proyek. & 0 \\
3 & SCPMK0802-01401 & Model Prototyping dan Spiral. & Bentuk: Ceramah, diskusi studi kasus, dan latihan analisis. Metode: Case Method. & 1 x (4 x 50') tatap muka; tugas terstruktur dan mandiri. & Mahasiswa mengkaji literatur model SDLC dan berdiskusi membandingkan model Prototyping dan Spiral. & Kriteria: ketepatan analisis kasus. Bentuk: dokumen analisis studi kasus (laporan kelompok). & Ketepatan menjelaskan karakteristik model beserta kelebihan dan kekurangannya. & 2.5 \\
4 & SCPMK0802-01401 & Model RUP (Rational Unified Process) dan analisis pemilihan metode. & Bentuk: Ceramah, diskusi studi kasus, dan latihan analisis. Metode: Case Method. Penugasan: analisis skenario proyek fiktif untuk menentukan metode Prototyping, Spiral, atau RUP yang tepat. & 1 x (4 x 50') tatap muka; tugas terstruktur dan mandiri. & Mahasiswa menganalisis kasus pemilihan metode dan menyusun dokumen analisis perbandingan beserta alasan pemilihan. & Kriteria: ketepatan analisis kasus. Bentuk: dokumen analisis studi kasus (laporan kelompok). & Ketepatan analisis risiko dan kejelasan kebutuhan; ketepatan pemilihan metodologi; logika argumentasi. & 2.5 \\
5 & SCPMK0802-01401 & Kuis 1: konsep dasar RPL, SDLC, dan metodologi preskriptif (Waterfall, Prototyping, Spiral, RUP). & Bentuk: Ujian tulis (daring/luring). Metode: evaluasi sumatif (closed book). & 1 x (2 x 50') ujian tulis. & Mahasiswa mengerjakan soal ujian secara mandiri, jujur, dan disiplin. & Kriteria: ketepatan jawaban sesuai kunci. Bentuk: tes objektif (pilihan ganda) dan/atau esai singkat. & Ketepatan menjelaskan tahapan SDLC; ketepatan membedakan karakteristik metodologi. & 2.5 \\
6 & SCPMK0802-01401 & Agile Methodology; pembentukan tim Scrum dan penyusunan Product Backlog. & Bentuk: Kuliah dan workshop (simulasi). Metode: Project Based Learning (PBL) fase inisiasi proyek. & 1 x (4 x 50') tatap muka; tugas terstruktur dan mandiri. & Mahasiswa membagi peran dalam tim (PO, Scrum Master) dan mengidentifikasi fitur aplikasi dalam bentuk User Stories. & Kriteria: kelengkapan artefak Scrum. Bentuk: penilaian kinerja proyek dan dokumen Product Backlog. & Ketepatan penulisan User Stories (INVEST); kejelasan prioritas Product Backlog. & 5 \\
7 & SCPMK0802-01401 & Scrum Framework dan simulasi Sprint Planning. & Bentuk: Kuliah dan workshop (simulasi). Metode: Project Based Learning (PBL) fase inisiasi proyek. & 1 x (4 x 50') tatap muka; tugas terstruktur dan mandiri. & Mahasiswa melakukan simulasi perencanaan Sprint 1 dalam tim Scrum. & Kriteria: kesesuaian prosedur events Scrum. Bentuk: penilaian kinerja proyek. & Kesesuaian pembagian peran tim; kejelasan definisi Done (DoD). & 5 \\
8 & SCPMK0802-01401, SCPMK0211-01403 & UTS: prinsip dasar RPL dan SDLC, metodologi preskriptif dan Agile, analisis kebutuhan (SRS), dan pemodelan visual sistem (UML). & Bentuk: Ujian tulis. Metode: evaluasi sumatif mandiri. & 1 x (2 x 50') ujian tulis. & Mahasiswa mengerjakan soal studi kasus analisis dan desain sistem secara mandiri. & Kriteria: kebenaran analisis dan ketepatan notasi desain. Bentuk: soal uraian (essay) dan studi kasus perancangan. & Ketepatan alasan pemilihan metodologi; ketepatan transformasi kebutuhan menjadi diagram UML yang konsisten. & 7.5 \\
9 & SCPMK0607-01402 & Metode pengumpulan kebutuhan (analisis dan elisitasi); dokumen SRS standar IEEE 830. & Bentuk: Kuliah dan studi kasus. Metode: Project Based Learning (tahap analisis). Penugasan: membuat dokumen SRS. & 1 x (4 x 50') tatap muka; tugas terstruktur dan mandiri. & Mahasiswa melakukan simulasi wawancara dengan klien, mengidentifikasi fitur, dan mendokumentasikannya ke dalam standar dokumen SRS (IEEE 830). & Kriteria: kelengkapan dan ketidakambiguan spesifikasi. Bentuk: dokumen SRS. & Ketepatan identifikasi kebutuhan fungsional/non-fungsional; kejelasan bahasa dalam SRS. & 10 \\
10 & SCPMK0211-01403 & Desain perangkat lunak dan pemodelan UML. & Bentuk: Workshop desain. Metode: Project Based Learning (tahap desain). Penugasan: membuat dokumen desain (SDD). & 1 x (4 x 50') tatap muka; tugas terstruktur dan mandiri. & Mahasiswa menerjemahkan kebutuhan (SRS) menjadi diagram visual (UML) menggunakan tools seperti StarUML/Draw.io. & Kriteria: konsistensi antar diagram. Bentuk: laporan desain sistem. & Ketepatan sintaks UML; logika alur pada Sequence Diagram; struktur database pada Class Diagram. & 5 \\
11 & SCPMK0602-01404 & Arsitektur perangkat lunak dan pola desain (Design Patterns). & Bentuk: Kuliah dan praktikum. Metode: Project Based Learning. Penugasan: implementasi pola desain pada kerangka proyek. & 1 x (4 x 50') tatap muka; tugas terstruktur dan mandiri. & Mahasiswa merancang struktur folder proyek berdasarkan arsitektur (misal MVC) dan menerapkan minimal satu Design Pattern dalam kode. & Kriteria: efisiensi struktur kode. Bentuk: project blueprint. & Ketepatan pemilihan arsitektur sesuai skala proyek; penerapan pola desain yang benar. & 10 \\
12 & SCPMK0802-01401 & Prinsip Kanban, papan visual (To Do, Doing, Done), dan Work In Progress (WIP) Limits. & Bentuk: Workshop manajemen proyek. Metode: Project Based Learning. Penugasan: setup project board (Trello/Jira/GitHub Projects). & 1 x (4 x 50') tatap muka; tugas terstruktur dan mandiri. & Mahasiswa mempraktikkan manajemen proyek dengan papan Kanban, memecah fitur menjadi kartu tugas, menetapkan WIP limits, dan melakukan daily stand-up. & Kriteria: kedisiplinan proses, transparansi alur kerja, dan kerjasama tim. Bentuk: observasi log aktivitas pada tools manajemen proyek. & Ketepatan pengelolaan task; kepatuhan alur kerja; kolaborasi dan pemerataan beban kerja. & 5 \\
13 & SCPMK0205-01405 & Implementasi (Construction/Coding). & Bentuk: Praktikum lab. Metode: Project Based Learning (tahap konstruksi). Penugasan: pengembangan aplikasi (coding). & 1 x (4 x 50') tatap muka; tugas terstruktur dan mandiri. & Mahasiswa melakukan coding secara kolaboratif menggunakan Git dan melakukan code review antar anggota tim. & Kriteria: aplikasi berjalan sesuai fungsi. Bentuk: source code dan demo aplikasi. & Modularitas kode; tidak ada error/bug kritikal; penerapan Clean Code. & 20 \\
14 & SCPMK0802-01401 & Teknik pengujian (Black Box vs White Box), penyusunan Test Case, dan pelaporan bug (Bug Tracking). & Bentuk: Workshop testing. Metode: Problem Based Learning. Penugasan: dokumen uji (Test Report). & 1 x (4 x 50') tatap muka; tugas terstruktur dan mandiri. & Mahasiswa membuat skenario uji (Test Case) untuk fitur tim lain (cross-testing) dan melaporkan temuan kesalahan. & Kriteria: akurasi temuan bug. Bentuk: laporan pengujian. & Kelengkapan skenario uji (positif dan negatif); kejelasan laporan bug. & 5 \\
15 & SCPMK0802-01401 & Evolusi perangkat lunak (Maintenance). & Bentuk: Kuliah dan presentasi akhir. Metode: diskusi dan refleksi proyek. Penugasan: user manual dan final report. & 1 x (4 x 50') tatap muka; tugas terstruktur dan mandiri. & Mahasiswa menganalisis potensi perubahan sistem di masa depan dan menyusun panduan penggunaan (User Manual). & Kriteria: kesiapan produk untuk rilis. Bentuk: presentasi proyek akhir. & Kelengkapan dokumentasi akhir; pemahaman siklus hidup software jangka panjang. & 0 \\
16 & SCPMK0205-01405, SCPMK0602-01404, SCPMK0802-01401 & UAS: arsitektur dan pola desain, implementasi kode (Clean Code), pengujian (Testing dan SQA), dan pemeliharaan (Maintenance). & Bentuk: Ujian akhir semester dan presentasi proyek. Metode: project presentation dan demo produk. Penugasan: submit laporan akhir dan demo aplikasi. & 1 x (3 x 50') presentasi dan demo. & Mahasiswa mempresentasikan hasil akhir produk perangkat lunak, mendemonstrasikan fitur, dan mempertahankan argumen teknis di hadapan penguji. & Kriteria: kualitas produk (bebas bug), kelengkapan fitur, dan kualitas arsitektur. Bentuk: rubrik penilaian proyek akhir dan ujian tulis (opsional). & Ketepatan penerapan arsitektur dan pola desain; kelengkapan skenario pengujian; kelengkapan dokumen maintenance. & 20 \\
17 & SCPMK0205-01405, SCPMK0602-01404, SCPMK0802-01401 & Pengayaan dan remedial seluruh materi pekan 1--16. & Bentuk: Pengayaan/remedial. Metode: diskusi individual dan feedback proyek. & Sesuai kebutuhan mahasiswa. & Mahasiswa yang belum mencapai Sub-CPMK mendapat bimbingan tambahan; mahasiswa yang sudah lulus melakukan pengayaan topik lanjutan. & Umpan balik formatif; tidak ada penilaian sumatif. & Pencapaian Sub-CPMK yang belum terpenuhi. & 0 \\
\end{longtblr}
